\uuid{ta3n}
\titre{Graphe}
\auteur{bodin ; hanani ; mzari}
\organisation{exo7}
\licence{CC BY-NC-SA 4.0}
\datecreate{2020-01-20}
\niveau{L1}
\module{Analyse}
\chapitre{Fonctions circulaires et hyperboliques inverses}
\sousChapitre{Autre}
\difficulte{2}
\theme{Fonctions usuelles, graphe}
\exo7qcmid{lille-1.fonctions-usuelles.017}

\contenu{
\begin{question}
\qclassification{126.00}
\qtags{motcle=graphe}
Soit $ f(x)= \frac{2x+1}{x-1}$. Quelles sont les assertions vraies ?
\begin{answers}
    \good{$y=2$ est une asymptote à la courbe de $f$ en $+\infty$.}

    \good{La courbe de $f$ admet une asymptote verticale $(x=1)$.}

    \bad{Le point de coordonnées $(1,1)$ est un centre de symétrie du graphe de $f$.}

    \good{Le point de coordonnées $(1,2)$ est un centre de symétrie du graphe de $f$.}
\end{answers}
\begin{explanations}
$\lim_{x\to +\infty}f(x)=2$, donc $y=2$ est une asymptote à la courbe de $f$ en $+\infty$.

$\lim_{x\to 1^+}f(x)=+\infty$ et $\lim_{x\to 1^-}f(x)=-\infty$, donc la droite $x=1$ est une asymptote (verticale) à la courbe de $f$ en $1$.

Le graphe de $f$ admet un centre de symétrie d'abscisse $1$ si et seulement si la fonction $x\to f(1+x)+f(1-x)$ est constante, pour tout $x\neq 0$. Ce qui revient à ce que la fonction $x\to f(x)+f(2-x)$ soit constante, pour tout $x\neq 1$. Or $f(x)+f(2-x)=4=2\times 2, \, \forall x \neq 1$. Donc le point de coordonnées $(1,2)$ est un centre de symétrie du graphe de $f$.
\end{explanations}
\end{question}
}
